% Research CV, Ekansh Bajpai. Build: pdflatex Research_CV.tex (twice), then copy Research_CV.pdf to ../pdfs/
\documentclass[10pt,a4paper]{article}
\usepackage[T1]{fontenc}
\usepackage[utf8]{inputenc}
\usepackage{charter}
\usepackage{microtype}
\usepackage{amssymb}
\usepackage{fontawesome5}
\usepackage[margin=1.2cm,top=0.8cm,bottom=0.6cm]{geometry}
\usepackage{xcolor}
\usepackage{enumitem}
\usepackage{titlesec}
\usepackage{tabularx}
\usepackage[hidelinks]{hyperref}

\definecolor{accent}{RGB}{31,78,120}
\newcommand{\link}[2]{\href{#1}{\color{accent}#2}}
\pagestyle{empty}
\setlength{\parindent}{0pt}
\setlist[itemize]{leftmargin=1.2em,itemsep=0pt,topsep=0pt,partopsep=0pt,parsep=0pt,
  label={\color{accent}\scriptsize$\blacktriangleright$}}
\titleformat{\section}{\large\scshape\color{accent}}{}{0pt}{}[\vspace{-0.3ex}{\color{accent}\titlerule}\vspace{0.2ex}]
\titlespacing*{\section}{0pt}{0.7ex}{0.4ex}

% \entry{title}{date or place}{subtitle line}
\newcommand{\entry}[3]{%
  \textbf{#1}\hfill{#2}\par
  \ifx&#3&\else{\small\itshape #3}\par\fi}

\begin{document}
\fontsize{9.8pt}{11.1pt}\selectfont

{\centering
{\Huge\scshape Ekansh Bajpai}\\[2pt]
{\large Robotics researcher: event-based vision, SLAM and robot learning}\\[3pt]
{\small \faMapMarker*\ Paderborn, Germany \quad
\faEnvelope\ \link{mailto:ekansh.bajpai712@gmail.com}{ekansh.bajpai712@gmail.com} \quad
\faGlobe\ \link{https://ekansh.is-a.dev}{ekansh.is-a.dev} \quad
\faGithub\ \link{https://github.com/ekansh-bajpai}{ekansh-bajpai} \quad
\faLinkedin\ \link{https://linkedin.com/in/ekanshbajpai}{ekanshbajpai}}\par}

\section{Research Interests}
Event-based and neuromorphic vision for visual odometry, SLAM and scene understanding; perception for robot manipulation
in dynamic scenes; learning-based robot control, from imitation learning with vision-language-action models to whole-body
control.

\section{Education}
\entry{M.Sc. Electrical Systems Engineering}{Paderborn University, Germany}{Thesis: Exploring Neuromorphic Front-Ends for Event-Based SLAM (\textbf{grade 1.0})}
\begin{itemize}
\item Focus: signal and information processing, SLAM, neural networks (CNNs, SNNs), event-based vision
\end{itemize}
\entry{B.Tech.\ Electronics and Communication Engineering}{2015 -- 2019}{Madan Mohan Malaviya University of Technology, India}
\begin{itemize}
\item Thesis: self-balancing object-tracking robot (PID-controlled inverted pendulum, IMU fusion, OpenCV tracking)
\end{itemize}
\textbf{IEEE CIS \& COLOURS Summer School}: \link{https://upb-rat.github.io/ieee-summer-school/}{\emph{Foundational Models Meet Robotics: From Language and Vision to Action}}\hfill Jul 2026\par
{\small\itshape Paderborn University, Germany}\par

\section{Research Experience}
\entry{Master's Thesis Research}{2026}{Automatic Control Group (RAT), Paderborn University}
{\small\itshape Supervisor: \link{mailto:alper.yegenoglu@uni-paderborn.de}{Dr.\ Alper Yegenoglu} $\cdot$ Examiners:
\link{mailto:erdal.kayacan@uni-paderborn.de}{Prof.\ Dr.\ Erdal Kayacan}, \link{mailto:sina.saravi@uni-paderborn.de}{Prof.\ Dr.\ Sina Saravi}}\par
\begin{itemize}
\item Replaced the CNN patch selector of the learned event-based VO pipeline DEVO with spiking neural networks (surrogate
      gradients); compared voxel-grid and time-surface encodings against CNN baselines in 150 runs
\item MVSEC: $\sim$\textbf{2$\times$ lower absolute trajectory error} and \textbf{10$\times$ smaller spread} than CNN baselines;
      held-out TartanEvents: both viable spiking variants beat both CNNs
\item Showed with a same-data CNN control that the gap on fast drone-racing motion (UZH-FPV) comes from training data, not
      spiking computation
\item Generated \textbf{TartanEvents}, a 1.3\,TB simulated event-camera dataset (open release after paper submission; access on request)
\item Built SuperEventSLAM, a graph-based SLAM back-end for DEVO (keyframes, SuperEvent features, three-stage loop
      verification, g2o); runs end to end at 36.7\,FPS, loop verification still being tuned
\end{itemize}
\entry{Master's Project \emph{Rescue Robots}: Autonomous Mapping and SLAM Optimization}{}{GET Lab (Prof.\ Dr.-Ing.\ Bärbel Mertsching), Paderborn University}
\begin{itemize}
\item LOAM-based lidar SLAM with IMU pose estimation; templated C++ architecture to benchmark pose-graph back-ends (g2o,
      SE-Sync) in ROS
\end{itemize}

\section{Work Experience}
\entry{Student Assistant (SHK), Humanoid Robotics}{Sep 2026 -- present}{Stratosfare -- Fraunhofer IEM, Paderborn}
\begin{itemize}
\item Implemented NVIDIA's GR00T end-to-end workflow on a real \textbf{Unitree G1} humanoid: teleoperated demonstrations,
      LeRobot dataset export, GR00T N1.7 VLA fine-tuning, LEAPP export and deployment on NVIDIA Jetson AGX Thor
\item Set up and integrated NVIDIA GR00T WholeBodyControl (GEAR-SONIC): robot interfaces, software and runtime environment;
      work directly on real humanoid hardware for teleoperation and whole-body motion control
\item Built the data-collection pipeline for learned whole-body control; collect whole-body VLA training data with PICO VR
      teleoperation
\item Demonstrated the trained G1 picking an object and placing it in a container at a 3-day workshop; guided attendees
      through teleoperation as used for episode collection
\end{itemize}
\begin{tabularx}{\linewidth}{@{}X@{\quad}r@{}}
\textbf{Student Assistant, AI/ML and Software}, Fraunhofer IEM: LLM-based knowledge and compliance platforms (RAG, FastAPI) & 2024 -- present \\
\textbf{Robotics Engineering Intern}, Cellgo GmbH: embedded I2C/SPI for a carrier robot & Aug 2022 \\
\end{tabularx}\par

\section{Teaching}
\entry{Student Assistant, Nonlinear Control of Autonomous and Robotic Systems}{Apr -- Sep 2026}{Automatic Control Group (RAT), Paderborn University}
\begin{itemize}
\item Built the ROS\,2/Gazebo quadrotor framework and tutorials (PID, feedback linearization, sliding mode, minimum-snap);
      supported Lyapunov stability analysis; graded term projects
\end{itemize}

\section{Publications, Datasets and Software}
\begin{itemize}
\item E.\ Bajpai. \emph{Exploring Neuromorphic Front-Ends for Event-Based SLAM}. Master's thesis, Paderborn University, 2026.
      Conference paper with my supervisor: in preparation.
\item \textbf{TartanEvents} (dataset, 1.3\,TB): \link{https://huggingface.co/datasets/upb-rat-ekansh/TartanEventsLeft_wo_upsampling}{huggingface.co/datasets/upb-rat-ekansh/TartanEventsLeft\_wo\_upsampling}
\item \textbf{devo-vo}: portable subset of DEVO for event-based VO (\link{https://github.com/ekansh-bajpai/devo-vo}{github.com/ekansh-bajpai/devo-vo});
      \textbf{pangolin-python-viz}: Python visualisation library on Pangolin for SLAM
      (\link{https://github.com/ekansh-bajpai/pangolin-python-viz}{github.com/ekansh-bajpai/pangolin-python-viz});
      snn-devo, supereventslam: private, access on request
\end{itemize}

\section{Skills}
\begin{tabularx}{\linewidth}{@{}l X@{}}
\textbf{Perception} & event cameras, spiking neural networks (snnTorch), visual odometry, SLAM, pose-graph optimization, bundle adjustment, OpenCV \\
\textbf{Robot learning} & imitation learning, VLA fine-tuning (GR00T, LeRobot), teleoperation, whole-body control, Unitree G1 \\
\textbf{Control} & nonlinear control, sliding mode, feedback linearization, Lyapunov stability, PID, trajectory generation \\
\textbf{Programming} & Python, C++, PyTorch, ROS\,2, Eigen, g2o, Gazebo, Git, Linux \\
\end{tabularx}

{\small\color{accent}\textbf{References:}} {\small \link{mailto:erdal.kayacan@uni-paderborn.de}{Prof.\ Dr.\ Erdal Kayacan} (Chair, RAT) $\cdot$
\link{mailto:alper.yegenoglu@uni-paderborn.de}{Dr.\ Alper Yegenoglu} (thesis supervisor, RAT), Paderborn University}
\end{document}
